\section{Related Work}\label{sec:related}
Previous works have focused on classifying different types of reward hacking and sometimes mitigating its effects. One popular setting is an agent on a grid-world with an erroneous sensor. \citet{hadfield2017ird} show and mitigate the reward hacking that arises due to an incorrect sensor reading at test time in a 10x10 navigation grid world. \citet{leike2017gridworld} show examples of reward hacking in a 3x3 boat race and a 5x7 tomato watering grid world. \citet{everitt2017corruptedreward} theoretically study and mitigate reward hacking caused by a faulty sensor.  

Game-playing agents have also been found to hack their reward. \citet{baker2020emergent} exhibit reward hacking in a hide-and-seek environment comprising 3-6 agents, 3-9 movable boxes and a few ramps: without a penalty for leaving the play area, the hiding agents learn to endlessly run from the seeking agents. \citet{toromanoff2019drl} briefly mention reward hacking in several Atari games (Elevator Action, Kangaroo, Bank Heist) where the agent loops in a sub-optimal trajectory that provides a repeated small reward. 

Agents optimizing a learned reward can also demonstrate reward hacking. \citet{ibarz2018humandemo} show an agent hacking a learned reward in Atari (Hero, Montezuma's Revenge, and Private Eye), where optimizing a frozen reward predictor eventually achieves high predicted score and low actual score. \citet{christiano2017preflearning} show an example of reward hacking in the Pong game where the agent learns to hit the ball back and forth instead of winning the point. \citet{stiennon2020learning} show that a policy which over-optimizes the learnt reward model for text summarization produces lower quality summarizations when judged by humans. 

